\section{Exact verification and reproducibility}
\label{sec:verification}

The infinite conclusions use the written proofs and the finite certificates
to which those proofs reduce the problem. Computations use integer or
rational arithmetic; no floating-point spectrum or Lean formalization is
claimed for these Laplacian results. SymPy~1.14.0 and Python~3.10 or later
are used, with source hashes in the saved results.

The \href{https://github.com/the-omega-institute/ideal-intersection-laplacian}{project repository}
contains every checker, exact certificate and reproduction command~\cite{repository}.
The build guide gives the commands by filename. The complete per-checker
catalogue is retained in \path{paper/sections/verification-details.tex};
the optional detailed build includes it without changing the proofs or
saved certificates. The principal verification scopes are as follows.

\subsection{Graph embeddings and repeated exponents}

Direct ideal adjacency validates $LU=UB$, independently of the support
construction, on $14$ squarefree and $(1,1,k)$ cases with $44\,294$
vertex pairs. The full $34$-vertex polynomial at $(2,2,3)$ is checked.
Three boundary graphs with $34,174,548$ vertices cover $165\,490$ pairs
and verify both repeated-exponent embeddings and the quotient lifting.
The $548$-vertex full characteristic polynomial is not computed.

The $90$ cubic diagnostics have exact rational intervals and independent
Sturm/sign checks. The all-$b$ proofs use the exhaustive divisor reduction,
not extrapolation from those diagnostics. All divisors of $24,108,320$
for smaller-factor indices $1,2,3$ are exhausted; eight positive-$b$
cubics have modular certificates checked by separate integer Horner
arithmetic. Higher indices use the uniform sign proof.
The separate ten-vertex graph for $(1,1,2)$ confirms why the symmetric
cubic alone is not an integrality characterization.

\subsection{Small minima and complete finite domains}

The general six-support quintic agrees with the reflected source quotient
and weighted symmetry. Thirty-five requested fully distinct triples and
direct ideal graphs with $22,58$ vertices check the construction. Endpoint
identities and positive shifted coefficients support the written unit,
minimum-two and uniform-cutoff proofs.

At minimum three, the proved cutoff leaves precisely $325$ pairs
$4\leq b<c<30$. All integer quintics, endpoint signs and intervals are
checked; the first endpoint has $257$ positive and $68$ negative values,
with no zero. This complete certificate is part of the minimum-three proof.
At minima four through seven, all $27\,562$ ordered triples in the derived
region are exhausted. A separate fraction-free Bareiss implementation
reconstructs every discriminant by a $9\times9$ integer Sylvester
determinant, checking exact divisions, consecutive-square brackets and
domain completeness. Together with the written cutoff, this proves
Corollary~\ref{cor:minimum-seven}; the minimum itself is not extrapolated.

\subsection{Inertia and endpoint surfaces}

Symbolic block, Schur-complement, determinant and positivity identities
support the uniform low root and balanced region. Six secular/Sturm
fixtures cover the possible nonsingular sign combinations; nine direct
quotient/Sturm fixtures include equal exponents and zero-diagonal boundaries.
The four bounded-gap families also have complete $38$-residue modular
certificates. These finite fixtures check formulas; the infinite conclusions
use the written comparison and inertia arguments.

For the middle-exponent bound, all $132$ cubics over $66$ pairs
$4\leq a<b\leq15$ are independently reconstructed. Integer trial division
exhausts $8\,527$ positive divisors, and separate Horner evaluation checks
all $7\,585$ above $b$. No upper bound on $c$ is imposed. This is the
complete certificate underlying Theorem~\ref{thm:endpoint-reduction}.
The linear cutoff uses the general endpoint cubic and five positive
coefficient expressions. Both known repeated-exponent endpoint zeros are
positive controls. The CRT scope note uses six substitution identities,
two positive boundary shifts and two integer seed determinants; its
all-moduli claim is the written polynomial-congruence argument.

\subsection{Arithmetic and shifted-root cases}

Scaling, all parity patterns and common residues for primes $3,5,7$
are checked coefficientwise, with $70$ quadratic-residue evaluations and
separate integer Horner arithmetic. The all-even scaling contradiction
uses the written valuation proof. Complete endpoint tables modulo two
and three include zero residues and zero polynomials; $70$ independent
integer determinant evaluations check the endpoint table.

For the cases in Appendix~\ref{sec:two-adic-lifting}, direct and reflected
quintics and all exponent permutations agree. Every shifted divisibility
and binary-cubic identity is checked coefficientwise for arbitrary shifts.
Chosen full quintics and derivatives are independently reconstructed by
integer determinants and principal cofactors, including compatible controls
and both known endpoint-zero examples. The finite residue and parity tables
support the identities; the infinite conclusions use the written root
factorizations, rather than a scan of exponents.

Reza's higher-modulus script was run unchanged; the independent verifier
checks all $28\,032$ residue triples with $56\,064$ Bareiss determinants.
The coefficientwise period-eight proof explains why the fixed divisor-$32$
tables add no classes at larger exponent moduli. His derivative script's
$384$ representatives are checked by $16\,380$ principal-minor determinants,
including the additional controls. Zero has infinite valuation, and full
valuations need not be constant on an exponent residue class. The complete
catalogue retains these diagnostic limits, coefficient counts, fixture
counts and exact source/certificate mappings.

\subsection{Endpoint-one spectral location}

Theorem~\ref{thm:endpoint-one-spectrum} and
Corollary~\ref{cor:endpoint-one-divisor-window} use written proofs without
a finite base. The
\href{https://github.com/the-omega-institute/ideal-intersection-laplacian/blob/d09cfdcc08f2a0b6c5b3abbd7d7a30bacc9030b4/scripts/check_endpoint_one_minimum_root.py}{minimum-root checker}
reconstructs the generic endpoint polynomial, comparison and equality
factorization, all twelve positive terms and the determinant identities.
The
\href{https://github.com/the-omega-institute/ideal-intersection-laplacian/blob/fd02ed726d3ada4bbd5e401e8818739f3bddc371/scripts/check_endpoint_one_pair_sum_window.py}{pair-sum checker}
verifies the four-support block, characteristic factor, discriminant and
strict sign identity. The established controls $(8,8,105)$, $(9,9,136)$
and $(20,20,741)$ leave $5,5,13$ divisors in the smallest-root windows;
all $23$ integer Horner evaluations are nonzero. Rational factorization
agrees, and exact Sturm counts give one quartic root in each window and
none on $[0,a]$. These controls already follow from the repeated-exponent
theorem. The source-hashed certificates reproduce byte-for-byte; their
finite diagnostics supplement the written bounds.

Proposition~\ref{prop:endpoint-one-equality-spectrum} also has a written
proof without a finite base. The
\href{https://github.com/the-omega-institute/ideal-intersection-laplacian/blob/d3fa3478193f34777a88b970da8f0763ce177c38/scripts/check_endpoint_one_equality_root.py}{equality-index checker}
verifies the principal factors, weighted symmetry and nonzero coupling.
The
\href{https://github.com/the-omega-institute/ideal-intersection-laplacian/blob/9e98266eb279a31b7282cea17fca3f5af6d9410f/scripts/check_endpoint_one_equality_second_root.py}{equality pair-sum checker}
verifies both positive identities at $a+b$. The three fixed controls
$(4,4,20)$, $(9,12,159)$ and $(9,30,951)$ are all off endpoint one.
After removing root $b$, exact Sturm counts place one positive root below
$a+b$ and three above it. These are diagnostics of the unconditional
statement, not integer endpoint solutions or a finite base.

\ifdefined\DetailedVerification
\subsection{Complete checker catalogue}
\begin{enumerate}
\item \path{scripts/check_integrality_obstructions.py} constructs ideals
as exponent tuples and tests adjacency by coordinatewise maxima, independently
of the support description. It verifies $LU=UB$ for the displayed
invariant embeddings, for squarefree $t=3,\ldots,8$ and $(1,1,k)$ with
$k=1,\ldots,8$: $14$ cases and $44\,294$ vertex pairs. It also checks
the square-gap inequalities and principal-minor witnesses.

\item \path{scripts/check_remaining_case.py} constructs all $34$ vertices
for $(2,2,3)$ and verifies its full characteristic polynomial and the
class decomposition in exact arithmetic. The shifted symmetric cubic is
$x^3-80x^2+2099x-18052$; its reduction modulo $5$ has no root.

\item \path{scripts/check_repeated_exponent_family.py} checks the general
quotient polynomial, both boundary sign identities, the positive-coefficient
expansion, and the integral antisymmetric eigenvalues symbolically.
Direct graphs for $(2,2,3),(3,3,10),(4,4,21)$ have $34,174,548$ vertices,
covering $165\,490$ vertex pairs. Exact neighbor counts verify both blocks
and their lifting. The check computes the seven-support quotient polynomial;
it does not compute the full $548\times548$ characteristic polynomial.

\item \path{scripts/check_aa_b.py} factors the $90$ diagnostic cubics
over $\mathbb Q$ and uses exact rational root isolation for a root between
consecutive integers. Independent Sturm counts and endpoint signs were
checked for every saved interval. Its modular search is restricted to
$2,3,5,7,11,13,17,19,23,29,31$; absence of a listed witness is not a proof
that no prime works for that cubic.

\item \path{scripts/check_fixed_repeated_exponents.py} reconstructs
the cubic from $B$, checks the discriminant, product, cutoff and reducible
factor identities, and enumerates every factor pair for $a=2,3,4$.
It verifies the admissibility conditions and all modular witnesses in
Tables~\ref{tab:divisors} and~\ref{tab:cubics}. The general congruence
and factor inequality are proved in Section~\ref{sec:bounds}.
\item \path{scripts/check_second_cutoff.py} checks the polynomial
remainder modulo $a+2$, enumerates every divisor of $24$, and verifies
the four root-free cubic reductions settling all $d=a+2$ cases.
It also checks the second-cutoff identities and the two nonsquare
discriminants completing $a=5,6$. This finite certificate follows the
proved divisor reduction; it is not a parameter-range search.
\item \path{scripts/check_factor_index_reduction.py} verifies the general
fixed-index divisibility and reconstruction identities, exhausts the
divisors of $K_2=108$, and certifies the unique positive-$b$ index-$2$
cubic modulo $13$. It checks the third-cutoff identities and the four
remaining discriminants for $a=7,8,9$. These finite checks exhaust the
proved candidate sets, with no arbitrary parameter cutoff.
\item \path{scripts/check_repeated_exponent_completion.py} reconstructs
the general cubic from the $4\times4$ quotient, checks both sign-interval
identities, and verifies the fourth-cutoff and strictly positive gap
identities. It exhausts every divisor of $24,108,320$ for indices
$1,2,3$, reproducing all eight positive-$b$ modular cubic certificates
with separate integer Horner arithmetic. The infinite remaining range is
settled by the sign proof of Theorem~\ref{thm:repeated}, not by a scan.
\item \path{scripts/check_six_support_quotient.py} verifies the general
six-support and complement quotients, their weighted symmetry and reflected
quintics, and the two determinant values used in Theorem~\ref{thm:unit}.
Exact rational factorization covers all thirty-five requested triples
$1\leq a<b<c\leq7$, with every matrix and quintic cross-checked against
\path{scripts/check_fully_distinct.py}. All saved rational intervals
pass independent Sturm counts and endpoint-sign checks.
Direct ideal adjacency verifies the quotient for
$(1,2,3)$ and $(2,3,4)$, using $22$ and $58$ vertices and $312$
representative-to-vertex pairs. The infinite unit-exponent conclusion
uses the written endpoint proof.
\item \path{scripts/check_minimum_two.py} verifies the complement
endpoint identities at $a=2$, the symmetric expression in $b+c$ and $bc$,
and both positive-coefficient expansions that exhaust the infinite range.
It reconstructs the single exceptional quintic at $(2,3,4)$ and checks
its rational endpoint values with separate Horner evaluation and a Sturm
count. No new parameter range is enumerated.
\item \path{scripts/check_distinct_tail.py} reconstructs the general
six-support complement quintic and checks its cubic first endpoint,
the uniform cutoff identity and positive remainder, and both
minimum-three second-endpoint expansions. The cutoff itself derives
the finite domain $4\leq b<c<30$, containing exactly $325$ pairs.
For every pair the checker reconstructs the integer quotient quintic,
matches its symbolic specialization, and checks the displayed root
interval using separate rational Horner arithmetic. Separate cubic
Horner evaluation confirms all first-endpoint values: $257$ positive,
$68$ negative and no zero. The $(3,4,5)$ rational exception also has a
Sturm count of one in $(1,3/2)$. The minimum-three theorem uses this
complete finite certificate; the general cutoff uses the written
positive-coefficient argument.
\item \path{scripts/check_low_spectrum.py} reconstructs the weighted
symmetric complement quotient and checks its block form, Schur complement,
determinant identity and ordered-parameter derivative. For the four
bounded-gap families it verifies the upper-inertia expressions and both
endpoint factorizations, positive shifted first-endpoint coefficients,
and all $38$ root-free modular residues with separate integer Horner
arithmetic. Eight chosen representatives, the four families at $a=4,20$,
have independently reconstructed integer quotient quintics and Sturm
counts of two roots in $(0,3)$; at $a=20$ both lie in $(2,3)$.
These examples check the formulas. The infinite families use the written
inertia proof and the complete modular exclusions, with no parameter scan.
\item \path{scripts/check_open_region2.py} runs the coauthor-requested
domain $4\leq a\leq7$, $a<b<c<4a^2-2a$. It constructs the generic
$H$ quotient determinant once and specializes its integer coefficients
for every one of the $27\,562$ triples. Every discriminant is nonsquare.
Original-quotient and complement endpoint signs are reported separately;
the general reflection identity is checked symbolically. Sixteen
representatives, including every first complement sign-pattern witness,
have separately reconstructed integer matrix characteristic polynomials,
complement endpoints and exact rational factorizations. The compatibility
entrypoint \path{scripts/check_open_region.py} runs the same corrected code.
\item \path{scripts/check_open_region_certificate.py} independently
reconstructs every one of the $27\,562$ discriminants as a $9\times9$
integer Sylvester determinant, using a separate fraction-free Bareiss
implementation. It checks every exact division, every strict consecutive-square
bracket and the complete ordered domain. This is the exhaustive finite
certificate used in Corollary~\ref{cor:minimum-seven}, combined with the
written cutoff, rather than extrapolation to arbitrary minimum exponent.
\item \path{scripts/check_mixed_inertia.py} verifies the comparison
determinant, negative Rayleigh value, equal-exponent case, positive
comparison gap and balanced-span identities symbolically in a rational
$3\times3$ congruence. Six exact secular/Sturm fixtures cover every
nonsingular diagonal-sign/secular-sign combination possible here.
Nine separately reconstructed quotient/Sturm fixtures include equal
exponents, two zero-diagonal boundaries, a strongly unbalanced triple
and chosen spans up to $1000$. The finite fixtures validate formulas;
Theorem~\ref{thm:uniform-low-root} and
Corollaries~\ref{cor:balanced-interval}--\ref{cor:balanced-span}
use the written comparison and positivity proofs, with no parameter scan.
\item \path{scripts/check_arithmetic_obstructions.py} verifies symbolic
integer quotient scaling, the equal-exponent quintic factorization and
its quadratic discriminant. It checks all eight parity patterns and
every nonzero common residue for primes $3,5,7$, with all $70$ quadratic
residue evaluations, including the binary factor, cross-checked by
integer Horner arithmetic. Six direct quotient fixtures verify scaling,
low-root and modular witnesses, and the existing boundary triple
$(9,9,136)$ checks that an integer endpoint is compatible with other
noninteger roots. The infinite conclusions use the written divisibility
and finite-field proofs, with no exponent scan.
\item \path{scripts/check_endpoint_reduction.py} reconstructs the general
endpoint-two cubic in the middle exponent and the two endpoint constant
terms in the largest exponent. It verifies both rational-cutoff identities
and five positive shifted-coefficient expressions. Three chosen even
gcd-two fixtures below the earlier quadratic cutoff have direct integer
quotient, independent Horner and exact Sturm checks. The existing repeated
triples $(9,9,136)$ and $(10,10,12)$ check the endpoint divisibility
conditions and retain nonlinear rational factors; the latter checks the
strict linear bound. The infinite bound and even-tail corollary use written
positivity and divisibility proofs, without an exponent scan or divisor-set
enumeration. The companion \path{scripts/check_endpoint_congruence_scope.py}
checks six substitution identities, two positive boundary shifts and the two
seed endpoint zeros by integer six-by-six determinants. Its all-moduli
conclusion uses the written polynomial-congruence argument; it does not
assert actual zeros or integer spectra for the constructed candidates.
\item \path{scripts/check_even_exponent_congruence.py} verifies integer
quotient scaling, the independent binary block characteristic polynomial,
and the coefficientwise congruences modulo two and thirty-two for the
scaled quintic. Three chosen gcd-two triples inside the preceding even
linear bound have separately reconstructed integer quotients, rational
factorizations and independent integer Horner values at one and three.
The infinite congruence-family theorem follows from the written divisibility
contradiction, without an exponent scan.
\item \path{scripts/check_endpoint_surfaces.py} is Reza's requested
diagnostic, corrected to use $|V(H)|$ in the complement reflection and
to include $b=15$. Exact rational linear-factor extraction finds no
endpoint roots for $4\leq a<b\leq15$, $c>b$.
\path{scripts/verify_endpoint_surfaces.py} independently reconstructs
the direct disjoint-support complement matrix and all $132$ endpoint
cubics over the $66$ pairs. It exhausts $8\,527$ positive constant-term
divisors by integer trial division and checks all $7\,585$ divisors
above $b$ with separate integer Horner arithmetic. Two previously
settled repeated-exponent positive controls recover their endpoint roots
and match full quotient factorizations. This complete finite certificate,
combined with the uniform low-root theorem and earlier cases, proves
Proposition~\ref{cor:middle-fifteen}. There is no upper cutoff on $c$ or
extension beyond the requested pair range.
\item \path{scripts/check_endpoint_residues.py} reconstructs the general
endpoint cubics and records all $26$ pair/endpoint rows for primes two
and three, including zero residues and zero polynomials. At all $35$
positive residue representatives, a separate $720$-term integer
determinant expansion checks both endpoint values, for $70$ checks.
The eight parity reductions and five simultaneous modulo-three exclusions
are verified against direct quotient characteristic polynomials.
Two selected fully distinct triples have independent endpoint determinants
and exact Sturm counts; two known repeated-exponent controls retain their
endpoint zeros. The infinite classes in
Proposition~\ref{thm:modulo-three-endpoints} use the written low-root proof.
This is a residue enumeration, not an exponent search.
\item \path{scripts/check_mixed_parity_congruence.py} checks the two
mixed binary reductions and four initial coefficient-divisibility identities.
It verifies all $558$ coefficients in twelve substitutions for the six
rows of Proposition~\ref{thm:mixed-parity-congruence}. Six positive residue
representatives and two selected fully distinct triples have $16$
independent $720$-term integer determinant checks at one and three,
matching separate integer Horner values. The fixtures retain nonlinear
rational factors. Complete lower-modulus tables give compatible linear
root multisets for all $12$ and $80$ unordered mixed-parity exponent
classes modulo four and eight, respectively; this compatibility is not
an integrality claim. The infinite theorem follows from the written
three-odd-root divisibility lemma.
\item \path{scripts/check_odd_root_distribution.py} verifies the generic
quintic's six exponent permutations and the coefficientwise identities
in Proposition~\ref{thm:odd-root-distribution}: $160$ coefficients for the
shifted polynomial's divisibility by eight, $160$ for its reduction
modulo two, $42$ endpoint coefficients modulo thirty-two and $51$
derivative coefficients modulo sixteen. Four selected fixtures have
$28$ independently expanded integer six-by-six determinants, including
six interpolation points and the zero-root check per fixture. A further
$24$ five-by-five principal cofactors independently reconstruct the
derivative at one. The reconstructed quintics and derivatives agree;
the infinite theorem uses the written odd-root distribution argument.
\item Reza's \path{scripts/check_2adic_higher_moduli.py} was run unchanged
at exponent moduli eight, sixteen and thirty-two. The independent
\path{scripts/verify_2adic_higher_moduli.py} checks all $28{,}032$ ordered
mixed-parity residue triples using $56{,}064$ integer Bareiss determinants
for the two endpoints. Twelve coefficientwise identities, with $560$
coefficients, establish coordinate period eight modulo thirty-two.
The complete $384$ base rows therefore encode the higher tables.
The endpoint-divisibility test excludes $36$, $288$ and $2304$ classes,
respectively. These are finite residue tables, and their survivors
do not assert integer endpoint zeros or integral spectra.
\item Reza's \path{scripts/check_higher_derivatives.py} was run unchanged
on its $384$ canonical representatives. Its interpretation was then
corrected: zero has valuation infinity, and full valuations need not be
constant on an exponent residue class. The independent
\path{scripts/verify_higher_derivatives.py} reconstructs all four values
at one by principal integer minors. Two new fixtures at three and four
lift controls at one bring the total to $16{,}380$ minor determinants.
Twelve coefficientwise identities determine the odd-root distribution
modulo four; $36$ further identities establish coordinate period eight
for the derivative residues at one and three modulo $16$, $8$ and $4$.
The second- and third-derivative thresholds add no exclusions with this
root information. A coefficientwise identity for $h_C(3)/16$ proves
Proposition~\ref{thm:odd-root-sum}; its $64$ role-ordered classes modulo
sixteen display $48$ exclusions, including $16$ beyond the earlier
divisor-$32$ test. Full joint representative patterns are saved separately.
\item \path{scripts/check_odd_root_lift.py} matches the direct complement
quintic to the reflected source and checks all six exponent permutations.
All $185$ coefficients of the second shifted polynomial are divisible
by $64$, and $185$ difference coefficients verify its binary cubic
identity for arbitrary integer shifts. Eight shift-parity patterns show
the irreducible quadratic in the excluded cases. Four fixtures, two
new exclusions and two compatible controls, independently reconstruct
their full quintics by $28$ integer determinant expansions, including
the zero-root checks. The $64$ additional role-ordered exclusions
modulo thirty-two follow from the written congruence and its class count;
no exponent or modulus range scan is needed for
Proposition~\ref{thm:odd-root-lift}.
\item \path{scripts/check_clustered_odd_roots.py} matches the direct
and reflected generic quintics, all exponent permutations, and the
exact factorization of $h_C(b)$. It checks $84$ difference coefficients
for the derivative congruence modulo sixty-four, and $249$ scaled
coefficients and $249$ binary-difference coefficients for each of the
two third-shift identities. Four derivative residues and sixteen binary
shift-parity patterns display the resulting necessary conditions.
Eight fixtures independently reconstruct full quintics with $56$ integer
six-by-six determinant expansions; another $48$ five-by-five principal
cofactors reconstruct their derivatives at $b$. Two fixtures have
$h_C(b)=0$ and satisfy the derivative divisor but fail the third-shift
condition. Four compatible controls retain nonlinear rational factors.
The arithmetic conditions in Proposition~\ref{thm:clustered-odd-roots} use
the root-clustering proof in Appendix~\ref{sec:two-adic-lifting}, with no
exponent or modulus range scan.
\item \path{scripts/check_mixed_endpoint_roots.py} matches the direct
and reflected quintics and all exponent permutations. For the three
mixed-parity patterns, coefficientwise checks verify the binary quintic,
divisibility of the shifted polynomial by eight, and its binary reduction
to $(y+1)^3$. A further $51$ coefficients give $h_C(2)\equiv4\pmod8$
for residue roles $(3,0,0)$. Eight full quintics are independently
reconstructed by $56$ integer six-by-six determinant expansions;
exact Sturm counts check the low-root fixtures. Two previous clustered-root
controls pass all prior necessary conditions but are excluded by the
linear tail. Both known endpoint-zero controls are retained.
Proposition~\ref{thm:one-odd-endpoints} and
Proposition~\ref{cor:mixed-linear-tail} follow from the written root-distribution
argument and the preceding low-root and endpoint-two proofs.
No exponent or modulus range is scanned.
\item The pinned PR9 \path{check_endpoint_two_middle_tail.py} verifies
the full endpoint Schur determinant, the diagonal factorization and
the middle-bound positive expansion, with five selected quotient/Horner/Sturm
fixtures and the repeated endpoint-zero control. Its unbounded statement
uses the written positive identity.
\item The pinned PR9 \path{check_endpoint_two_sharp_maximum.py} verifies
the inverse interval substitution, both nonnegative expressions and every
coefficient of the $84$-term residual in
Lemma~\ref{lem:endpoint-two-rectangle}. All $16$ residual vectors agree
with Appendix~\ref{sec:endpoint-two-identities}; eight direct fixtures
and two below-cutoff controls supplement the written proof.
\item At PR9, \path{check_endpoint_two_completion.py} checks the
$44$- and $112$-term identities and all $32$ coefficient vectors.
It exhausts exactly $8\,658$ necessary fully distinct triples, with
independent integer-Horner and $6\times6$ Bareiss evaluations in every
case. The complete CSV and JSON reproduce byte-for-byte. Seven direct
quotient/Horner/Sturm fixtures and the repeated control supplement the
proof; earlier minimum-through-seven certificates remain dependencies.
\item At PR9, \path{check_endpoint_one_equality_cubic.py} verifies
the explicit cubic's generic quotient identity and discriminant by an
independent $5\times5$ Sylvester determinant. Its $61$-term expansion,
divisor/quadratic identity and abstract square-discriminant counterexample
are exact checks, with no exponent scan or point enumeration.
\end{enumerate}

The symbolic and characteristic-polynomial checks use SymPy~1.14.0;
the first checker uses the Python standard library.
Saved results contain source hashes. The $10$-vertex polynomial in the
preceding remark was independently checked by a direct ideal-graph
construction; it is also reproducible by the command in the manuscript
build guide. No floating-point eigenvalue calculation is used as
integrality evidence.

\fi
